\chapter{CONCLUSIONS}
\label{ch:conclusions}

This laboratory activity successfully demonstrated the process of exploring and evaluating different UML drawing tools through structured teamwork and collaborative platforms. The experiment helped students understand how software tools can be researched, compared, and selected based on usability, features, and project requirements.

The use of Trello for task management and Discord for communication provided an organized workflow for assigning responsibilities, tracking progress, and sharing updates. These tools improved coordination among team members and ensured that all tasks were completed within the given timeframe.

Through individual research and group discussion, the team was able to compare multiple UML tools and identify Draw.io as the most suitable option for academic and collaborative use due to its simplicity, accessibility, and free availability. The activity also reinforced the importance of documentation, communication, and evidence-based tool selection in software engineering projects.

Overall, this laboratory session established a strong foundation for future Object-Oriented Analysis and Design activities by introducing collaborative workflows, UML tool evaluation, and structured reporting practices that are essential in real-world software development environments.
